% Diagrama de topología HaLow Mesh con Raspberry Pi 4 como clientes STA
% Muestra 3 Alfa Networks Tube AHM (MT7628) formando mesh IEEE 802.11s
% y 3 Raspberry Pi 4 conectados como STAs vía HaLow

\begin{figure}[H]
\centering
\begin{tikzpicture}[
    % Estilos de nodos
    meshgate/.style={rectangle, draw=blue!60, fill=blue!20, thick, minimum width=2.5cm, minimum height=1.2cm, rounded corners, font=\small\bfseries},
    meshpoint/.style={rectangle, draw=green!60, fill=green!20, thick, minimum width=2.5cm, minimum height=1.2cm, rounded corners, font=\small\bfseries},
    gateway/.style={rectangle, draw=orange!60, fill=orange!20, thick, minimum width=2.8cm, minimum height=1.4cm, rounded corners, font=\small},
    thread/.style={circle, draw=purple!60, fill=purple!10, thick, minimum size=0.8cm, font=\tiny},
    % Estilos de conexiones
    mesh/.style={<->, thick, draw=green!70!black, dashed},
    halow/.style={->, thick, draw=blue!70!black, dotted},
    ethernet/.style={-, thick, draw=black},
    threadlink/.style={-, thick, draw=purple!70!black, dotted}
]

% ============ CAPA 1: RED HALOW MESH (802.11s) ============

% Mesh Gate (Tube-01) - Root con backhaul
\node[meshgate] (tube01) at (0,6) {
    \textbf{Tube-01}\\
    \textcolor{blue}{\footnotesize Mesh Gate}
};
\node[below=0.1cm of tube01, font=\tiny, text width=3cm, align=center] {
    Alfa Tube AHM\\MT7628 + MM6108\\915 MHz @ 4 MHz BW\\+20 dBm TX
};

% Mesh Point 2 (Tube-02) - Relay intermedio
\node[meshpoint] (tube02) at (-4,3) {
    \textbf{Tube-02}\\
    \textcolor{green!50!black}{\footnotesize Mesh Point}
};
\node[below=0.1cm of tube02, font=\tiny, text width=3cm, align=center] {
    Relay 280m\\RSSI -72 dBm\\Airtime cost 512
};

% Mesh Point 3 (Tube-03) - Relay extremo
\node[meshpoint] (tube03) at (4,3) {
    \textbf{Tube-03}\\
    \textcolor{green!50!black}{\footnotesize Mesh Point}
};
\node[below=0.1cm of tube03, font=\tiny, text width=3cm, align=center] {
    Relay 520m\\RSSI -88 dBm\\Airtime cost 2048
};

% Enlaces mesh 802.11s (peer links bidireccionales)
\draw[mesh] (tube01) -- (tube02) node[midway, above, sloped, font=\tiny] {Peer Link 1};
\draw[mesh] (tube01) -- (tube03) node[midway, above, sloped, font=\tiny] {Peer Link 2};
\draw[mesh] (tube02) -- (tube03) node[midway, above, font=\tiny] {Peer Link 3};

% Etiqueta red mesh
\node[above=0.3cm of tube01, font=\small\bfseries, text=blue!70!black] {
    Red HaLow Mesh IEEE 802.11s (902-928 MHz)
};

% ============ CAPA 2: GATEWAYS RASPBERRY PI 4 (CLIENTES STA) ============
% SOLO 2 GATEWAYS: RPi4-01 conectado a Tube-02, RPi4-02 conectado a Tube-03

% Gateway 1 (RPi4-01) - Conectado a Mesh Point 2 (Tube-02)
\node[gateway] (rpi01) at (-4,0) {
    \textbf{RPi4-01}\\
    \footnotesize Gateway Edge
};
\node[below=0.1cm of rpi01, font=\tiny, text width=3.2cm, align=center] {
    HaLow STA (wireless)\\ThingsBoard Edge\\OTBR + Docker stack
};

% Gateway 2 (RPi4-02) - Conectado a Mesh Point 3 (Tube-03)
\node[gateway] (rpi02) at (4,0) {
    \textbf{RPi4-02}\\
    \footnotesize Gateway Edge
};
\node[below=0.1cm of rpi02, font=\tiny, text width=3.2cm, align=center] {
    HaLow STA (wireless)\\ThingsBoard Edge\\OTBR + Docker stack
};

% Conexiones HaLow (STAs hacia APs mesh) - SOLO a Mesh Points, NO al Mesh Gate
\draw[halow] (rpi01) -- (tube02) node[midway, right, font=\tiny] {HaLow STA};
\draw[halow] (rpi02) -- (tube03) node[midway, right, font=\tiny] {HaLow STA};

% Etiqueta capa gateways
\node[left=0.5cm of rpi01, font=\small\bfseries, text=orange!70!black, rotate=90, anchor=south] {
    2 Gateways Edge (HaLow STAs)
};

% ============ CAPA 3: NODOS THREAD (ejemplo) ============

% Nodos Thread conectados a RPi4-01 (via Mesh Point 2)
\node[thread] (t01) at (-4.8,-1.5) {T1};
\node[thread] (t02) at (-4,-1.5) {T2};
\node[thread] (t03) at (-3.2,-1.5) {T3};

% Nodos Thread conectados a RPi4-02 (via Mesh Point 3)
\node[thread] (t04) at (3.2,-1.5) {T4};
\node[thread] (t05) at (4,-1.5) {T5};
\node[thread] (t06) at (4.8,-1.5) {T6};

% Enlaces Thread (IEEE 802.15.4 mesh)
\draw[threadlink] (rpi01) -- (t01);
\draw[threadlink] (rpi01) -- (t02);
\draw[threadlink] (rpi01) -- (t03);
\draw[threadlink] (rpi02) -- (t04);
\draw[threadlink] (rpi02) -- (t05);
\draw[threadlink] (rpi02) -- (t06);

% Etiqueta capa Thread
\node[below=0.3cm of t05, font=\small\bfseries, text=purple!70!black] {
    Nodos Thread (IEEE 802.15.4 @ 2.4 GHz) - 30 total
};

% ============ BACKHAUL (ETHERNET LOCAL) ============

% Switch Ethernet local (arriba del Mesh Gate)
\node[rectangle, draw=black, fill=gray!10, thick, minimum width=2.5cm, minimum height=0.8cm] (switch) at (0,8) {
    \textbf{Switch Ethernet}\\
    \tiny Local network
};

% Conexión Ethernet desde Mesh Gate a Switch
\draw[ethernet] (tube01) -- (switch) node[midway, right, font=\tiny] {GbE};

% Nota sobre server on-premise
\node[above=0.1cm of switch, font=\tiny, text width=4cm, align=center, text=gray!70!black] {
    Server on-premise conectado\\al mismo switch (Cap. 5)
};

% ============ LEYENDA ============

\node[below=2cm of t02, font=\footnotesize, text width=5cm, align=left] (legend) {
    \textbf{Leyenda:}\\[0.2cm]
    \tikz\draw[mesh] (0,0) -- (0.5,0); Peer Link 802.11s\\
    \tikz\draw[halow] (0,0) -- (0.5,0); HaLow STA to AP\\
    \tikz\draw[threadlink] (0,0) -- (0.5,0); Thread 802.15.4\\
    \tikz\draw[ethernet] (0,0) -- (0.5,0); Ethernet GbE
};

\end{tikzpicture}
\caption{Topología HaLow Mesh con Gateways Raspberry Pi 4 como clientes STA. Red HaLow (capa superior) comprende 3 routers Alfa Networks Tube AHM formando mesh IEEE 802.11s con protocolo HWMP: 1 Mesh Gate (Tube-01) con backhaul Ethernet a switch local (server on-premise conectado al mismo switch, detallado Capítulo 5), y 2 Mesh Points (Tube-02, Tube-03) relay. Gateways Raspberry Pi 4 (capa intermedia, 2 unidades) se conectan como clientes STA inalámbricos exclusivamente a los Mesh Points (NO al Mesh Gate), ejecutando ThingsBoard Edge y OTBR. Nodos Thread (capa inferior) forman mesh IEEE 802.15.4 conectado a cada gateway vía nRF52840 Dongle USB. Este capítulo enfoca en red HaLow mesh 802.11s y edge gateways; conectividad WAN se documenta Capítulo 5.}
\label{fig:halow-mesh-topology}
\end{figure}
